\chapter{Expert Interview Protocol and Summary}\label{app:interview}
% Nguon: 03_phan-tich/phong-van/2026-07-09_interview-protocol_v3.0_draft_AI_EN.md (protocol)
%        03_phan-tich/phong-van/2026-07-10_phan-tich-phong-van_v1.0_draft_AI.md (CG-01, CG-02)
%        03_phan-tich/phong-van/expert_result_03.md, expert_result_04.md, expert_result_05.md

\section{Interview Protocol}\label{app:interview-protocol}

The semi-structured interview protocol used with all five experts was constructed
following the Interview Protocol Refinement (IPR) framework
\parencite{castillo-montoya_preparing_2016}, as summarised in
\Cref{subsec:expert-interviews}. The protocol comprised thirteen core questions organised
into the five parts listed in \Cref{tab:protocol-structure}, preceded by an informed-consent
script explaining the academic purpose of the study, the anonymisation of personal
information, and the participant's right to decline any question or withdraw at any time.
Each session was expected to last twenty to thirty minutes and was recorded through written
notes only, without audio recording.

\begin{table}[htbp]
  \centering
  \caption{Structure of the five-part interview protocol.}
  \label{tab:protocol-structure}
  \begin{tabular}{@{}p{2.6cm}p{5.2cm}p{6cm}@{}}
    \toprule
    \textbf{Part} & \textbf{Content} & \textbf{Purpose} \\
    \midrule
    0 & Greeting and informed consent & Establish ethical basis and obtain verbal consent \\
    1 --- Warm-up (Q1--Q3) & Recent, concrete laptop-selection episode & Ground the respondent's perspective in a specific case rather than an abstract opinion \\
    2 --- Stimulus-based critique (Q4--Q7) & A literature finding for each criterion group is read aloud immediately before the matching question & Anchor agreement or disagreement to a stated academic claim \\
    3 --- Vietnamese context (Q8--Q9) & Open questions on locally specific factors & Surface criteria the literature synthesis could not have anticipated \\
    4 --- Trade-off scenario (Q10--Q11) & Hypothetical comparisons between conflicting criteria & Generate qualitative seed data for the AHP pairwise comparisons of \Cref{sec:stage2} \\
    5 --- Closing (Q12--Q13) & Member-checking and open addition & Confirm the interviewer's understanding and capture unprompted remarks \\
    \bottomrule
  \end{tabular}
  \floatnote{\emph{Source:} authors' own protocol design, following
  \textcite{castillo-montoya_preparing_2016} and \textcite{willis_cognitive_2005}.}
\end{table}

Part 2 read a short stimulus card to the expert before each question, presenting an
academic finding as material for critique rather than as a leading prompt: a card on
hardware and performance reported that CPU and RAM were the two criteria rated most
important by Maghsoudi et al.\ (2026), with GPU ranking tenth; a card on mobility and
durability reported that Lam et al.\ (2023) ranked battery life as the most important
technical-specification criterion for mobile devices; and a card on experience and trust
reported that Sönmez Çakır and Pekkaya (2020) identified brand image as the most important
criterion overall, alongside the finding of Maghsoudi et al.\ (2026) that fan noise and
charger bulk drew frequent complaints yet had little influence on the final purchase
decision. Most substantive questions carried a follow-up probe asking the respondent to
recall the specific case behind an answer or to explain the reasoning process that produced
it, a technique adapted from cognitive interviewing \parencite{willis_cognitive_2005}.

\section{Expert Profiles}\label{app:interview-profiles}

Five experts were selected purposively so that each occupies a different position in the
laptop procurement process, as summarised in \Cref{tab:expert-profiles}.

\begin{table}[htbp]
  \centering
  \caption{Profiles of the five interviewed experts.}
  \label{tab:expert-profiles}
  \begin{tabular}{@{}lp{4.2cm}cp{3.6cm}c@{}}
    \toprule
    \textbf{Code} & \textbf{Role} & \textbf{Experience} & \textbf{Position in procurement process} & \textbf{Date} \\
    \midrule
    CG-01 & Deputy head, international relations office & 7 years & End-user with direct purchasing experience & 2026-07-05 \\
    CG-02 & IT administrator & 3 years & Specifies and maintains office device fleets & 2026-07-09 \\
    CG-03 & Import--export documentation officer & 3.5 years & End-user representing the target persona & 2026-07-10 \\
    CG-04 & Head of product development & 10 years & Manager with organisational purchasing authority & 2026-07-08 \\
    CG-05 & Retail sales and technical-support consultant & 3 years & Advises laptop buyers daily at point of sale & 2026-07-10 \\
    \bottomrule
  \end{tabular}
  \floatnote{\emph{Source:} authors' own interview data.}
\end{table}

\section{Thematic Summary of Findings}\label{app:interview-summary}

Interview notes were analysed thematically, cross-referencing each recurring position
against the two or more experts who independently raised it, consistent with the majority
corroboration rule described in \Cref{subsec:expert-interviews}.

\subsection{Hardware Priorities and the GPU Question}

Four of the five experts converged independently on removing the dedicated graphics card
from the criteria set. CG-02 characterised a discrete GPU as over-specification for office
work that adds cost, weight, heat, and battery drain without a matching benefit, a position
consistent with the tenth-place GPU ranking reported by \textcite{maghsoudi_hybrid_2026}
and read to the experts as Card A. CG-03, CG-04, and CG-05 reached the same conclusion from
their respective vantage points as a documentation officer, a product-development manager,
and a retail consultant, each reporting that office customers who ask about GPU capability
are rare and that the budget it would consume is better redirected to the CPU, RAM, or SSD.
CG-05 qualified the point by noting that a discrete GPU remains justified only when the
organisation has an occasional design or video-editing need, a caveat also raised by CG-02.
CG-01 was the sole dissenting voice: in the direct question on GPU priority, CG-01 described
integrated graphics as adequate for most office machines but characterised a discrete GPU as
``nice to have, no harm in not having it,'' and in the Part~4 trade-off scenario explicitly
chose the GPU-equipped option, reasoning that no buyer can be certain in advance that a
device will never be used for design or gaming work, so a more versatile machine is the
safer choice. This single dissent is the reason the criteria revision rule in
\Cref{subsec:expert-interviews} requires majority rather than unanimous corroboration before
a candidate criterion is changed. On RAM, all five experts converged on 16\,GB as the new
practical minimum, citing browser multitasking and simultaneous use of communication
platforms such as Microsoft Teams and Zalo (CG-02, CG-03) as the drivers of memory pressure
that an 8\,GB configuration can no longer absorb.

\subsection{Battery, Weight, and Design Trade-offs}

Four of the five experts (CG-02 through CG-05) explicitly attributed a lower priority on
maximum battery capacity to the ready availability of mains power in Vietnamese offices,
qualifying the emphasis placed on battery life by \textcite{weng_siew_lam_evaluation_2023},
read to the experts as Card B; CG-01 reached the same ranking of priorities without citing
mains-power access as the reason. CG-01,
CG-03, and CG-05 stated that a runtime of roughly half a working day, sufficient to cover a
meeting or a short period away from a desk, is adequate, while CG-04 set a more concrete
range of four to seven hours. In its place, physical weight emerged as the more decisive
mobility criterion: CG-01, CG-03, CG-04, and CG-05 all linked this preference to motorbike
commuting and daily carrying, converging on a practical range of roughly 1.3 to 1.6\,kg for
a 14-inch machine. On design, CG-01, CG-03, and CG-05 associated a slim, neutral-coloured
form factor with a professional impression in client-facing work, but CG-02 and CG-04
cautioned that this aesthetic preference is bounded by durability requirements, specifically
hinge rigidity, thermal headroom, and resistance to chassis flex, so that appearance is
treated as a secondary criterion once a minimum durability bar is met.

\subsection{Brand, After-Sales, and Divergent Views on Trust}

Interview responses on brand and after-sales service were the least uniform of the five
themes. CG-01 argued that the growing accessibility of AI-assisted comparison tools is
reducing buyers' reliance on brand reputation as a proxy for quality, a view echoed by
CG-03's observation that end-users increasingly research specifications independently.
CG-02, CG-04, and CG-05, however, each reported from a business-procurement or retail
vantage point that brand and after-sales reliability remain decisive when device downtime
carries an organisational cost: CG-02 emphasised that a business buyer facing two
comparably priced machines will choose the more established brand unless an internal IT
team can independently verify and service an unfamiliar one, while CG-04 and CG-05 both
noted that buyers who cannot evaluate technical specifications in depth default to two or
three familiar brands as a risk-management heuristic. CG-05 added that proximity of an
authorised service centre to the workplace is frequently a deciding factor for
organisational buyers, a point later adopted as the operational rationale for measuring
Brand through a service-centre-count proxy in \Cref{subsec:operationalisation}.

\subsection{Locally Specific Factors}

All five experts identified factors specific to the Vietnamese working environment that the
international literature reviewed in \Cref{sec:prior-studies} does not address directly.
The hot and humid climate was raised by CG-02 and CG-04 as a driver of thermal-management
and build-quality requirements beyond what temperate-climate studies typically consider.
Motorbike commuting was raised independently by CG-01, CG-03, CG-04, and CG-05 as the
mechanism linking weight and durability to daily use rather than occasional travel. CG-01
and CG-03 both described a software lock-in effect arising from Windows-only accounting,
customs-declaration, and internal enterprise-resource-planning software, which excludes
non-Windows devices from consideration regardless of hardware merit; this observation
directly motivated the addition of software compatibility as a replacement criterion for
GPU, as reported in \Cref{subsec:interview-outcomes}. CG-02 and CG-05 further noted that
the residual prevalence of HDMI projectors and USB Type-A peripherals in Vietnamese offices
keeps traditional ports relevant even as ultraportable designs move toward USB-C only.

\subsection{Trade-off Scenario Responses}

When asked in Part 4 to name the three criteria they would prioritise if all others were
kept flexible, core hardware performance (CPU, RAM, SSD, or a combination) appeared in every
one of the five answers, making it the only criterion named by all experts. Brand or
after-sales support was named by four of the five (CG-02 through CG-05), and durability was
named separately by three (CG-02, CG-03, CG-04). Price, although established earlier in the
interview as the first screening criterion applied by all five experts, recurred explicitly
in only two of the five top-three answers (CG-01, CG-05), suggesting that once a price
bracket has already narrowed the option set, experts treat it as a pre-filter rather than a
continuing point of comparison. This pattern of convergence on hardware performance as the
dominant criterion, together with the comparatively low priority assigned to battery life
and to size or weight once the office context was made explicit, foreshadows the weight
concentration later confirmed quantitatively by the AHP results in \Cref{sec:weights}.
